\begin{table*}[tp]
\centering
\footnotesize
\setlength{\tabcolsep}{2.5pt}
\renewcommand{\arraystretch}{1.1}
\renewcommand{\thetable}{S1}
\caption{Rendered-RGB triangulation reference}
\label{tab:s1}
\begin{tabular}{lcccccc}
\toprule
Construction / map & Landmarks & Map size (MB) & Median TE / RE & Mean TE (cm) & R5 (\%) & R15 (\%) \\
\midrule
Rendered RGB tracks (full) & 553,339 & 577.69 & 8.772 / 0.126 & 28.80 & 43.58 & 72.01 \\
AnyGSLoc candidates (full) & 535,840 & 559.43 & 8.762 / 0.127 & 42.49 & 43.00 & 71.99 \\
Rendered RGB tracks + spatial sampling & 176,329 & 184.10 & 13.785 / 0.181 & 71.94 & 31.42 & 59.18 \\
AnyGSLoc candidates + spatial sampling & 176,329 & 184.10 & 13.790 / 0.183 & 69.95 & 31.59 & 59.81 \\
AnyGSLoc compact & 175,346 & 183.07 & 8.812 / 0.129 & 38.64 & 42.38 & 69.79 \\
\bottomrule
\end{tabular}
\par\smallskip\begin{minipage}{0.98\linewidth}\footnotesize The rendered-RGB reference uses feature tracks, fixed-camera triangulation, and ordinary descriptor averaging; it is a self-implemented SfM-style reference. Both full maps are unselected. Spatial rows retain 176,329 points; our current compact map retains 175,346. This comparison changes map construction and is separate from the shared-observation geometry control.\end{minipage}
\end{table*}
